% Exercitium2013.tex
% Exercitium (curso 2013-14)
% José A. Alonso Jiménez <jalonso@us.es>
% Sevilla, 9 de diciembre de 2018
% ======================================================================

\documentclass[a4paper,12pt,twoside]{book}

%%%%%%%%%%%%%%%%%%%%%%%%%%%%%%%%%%%%%%%%%%%%%%%%%%%%%%%%%%%%%%%%%%%%%%%%
%% § Paquetes adicionales
%%%%%%%%%%%%%%%%%%%%%%%%%%%%%%%%%%%%%%%%%%%%%%%%%%%%%%%%%%%%%%%%%%%%%%%%

% Configuración para XeLaTeX
\usepackage{fontspec}
\usepackage{xltxtra}
\defaultfontfeatures{Ligatures=TeX,Numbers=OldStyle}
\setromanfont{DejaVu Sans}
% \setsansfont{Arial}
\setmonofont{DejaVu Sans Mono}[Scale={0.90}]

% Notas: La lista de fuentes disponibles se obtiene con fc-list

% \usepackage{ucs}
% \usepackage[utf8]{inputenc}        % Acentos de UTF8
\usepackage[spanish]{babel}        % Castellanización.
% \usepackage[T1]{fontenc}           % Codificación T1 con European Computer
%                                    % Modern.  
% \usepackage{graphicx}
% \usepackage{fancyvrb}              % Verbatim extendido
% \usepackage{mathpazo}              % Fuentes semejante a palatino
% \usepackage[scaled=.90]{helvet}
% \usepackage{cmtt}
% \renewcommand{\ttdefault}{cmtt}
\usepackage{a4wide}
\usepackage{minted}
\usepackage{comment}

\usepackage{titletoc}
\dottedcontents{chapter}[0em]{}{9em}{1pc}
% \dottedcontents{chapter}[<left>]{<above-code>}{<label width>}{<leader width>}

\linespread{1.05}                  % Distancia entre líneas
\setlength{\parindent}{2em}        % Indentación de comienzo de párrafo
\setlength{\parskip}{1ex}          % Distancia entre párrafos
% \deactivatetilden                  % Elima uso de ~ para la eñe
\raggedbottom                      % No ajusta los espacios verticales.

\usepackage[%
  colorlinks=true,
  urlcolor=blue,
  % pdftex,
  pdfauthor={José A. Alonso <jalonso@us.es>},%
  pdftitle={Exercitium (curso 2013-14)},%
  pdfstartview=FitH,%
  bookmarks=false]{hyperref}      

\setcounter{tocdepth}{1}
\setcounter{secnumdepth}{4}

\usepackage{tocstyle}
\usetocstyle{KOMAlike}

% \usepackage{tocloft}
% \renewcommand\cftpartnumwidth{3cm}

% \usepackage{minitoc}

% \setlength\cftparskip{-2pt}
% \setlength\cftbeforechapskip{0pt}

%%%%%%%%%%%%%%%%%%%%%%%%%%%%%%%%%%%%%%%%%%%%%%%%%%%%%%%%%%%%%%%%%%%%%%%%%%%%%%
%% § Cabeceras                                                              %%
%%%%%%%%%%%%%%%%%%%%%%%%%%%%%%%%%%%%%%%%%%%%%%%%%%%%%%%%%%%%%%%%%%%%%%%%%%%%%%

\usepackage{fancyhdr}

\addtolength{\headheight}{\baselineskip}

\pagestyle{fancy}

\cfoot{}

\fancyhead{}
\fancyhead[RE]{\mdseries\sffamily Exercitium (2013--14)}
\fancyhead[LO]{\mdseries\sffamily \nouppercase{\leftmark}}
\fancyhead[LE,RO]{\mdseries\sffamily \thepage}

%%%%%%%%%%%%%%%%%%%%%%%%%%%%%%%%%%%%%%%%%%%%%%%%%%%%%%%%%%%%%%%%%%%%%%%%
%% § Definiciones
%%%%%%%%%%%%%%%%%%%%%%%%%%%%%%%%%%%%%%%%%%%%%%%%%%%%%%%%%%%%%%%%%%%%%%%%

\input definiciones
\def\mtctitle{Contenido}

%%%%%%%%%%%%%%%%%%%%%%%%%%%%%%%%%%%%%%%%%%%%%%%%%%%%%%%%%%%%%%%%%%%%%%%%
%% § Título
%%%%%%%%%%%%%%%%%%%%%%%%%%%%%%%%%%%%%%%%%%%%%%%%%%%%%%%%%%%%%%%%%%%%%%%%

\title{
  {\LARGE Exercitium (curso 2013--14) \\
  {\Large Ejercicios de programación funcional con Haskell }}}
\author{\href{http://www.cs.us.es/~jalonso}
        {\Large José A. Alonso Jiménez}}
\date{\vfill \hrule \vspace*{2mm}
  \begin{tabular}{l}
      \href{http://www.cs.us.es/glc}
           {Grupo de Lógica Computacional} \\
      \href{http://www.cs.us.es}
           {Dpto. de Ciencias de la Computación e Inteligencia Artificial} \\
      \href{http://www.us.es}
           {Universidad de Sevilla}  \\
      Sevilla, 9 de diciembre de 2018
  \end{tabular}\hfill\mbox{}}

%%%%%%%%%%%%%%%%%%%%%%%%%%%%%%%%%%%%%%%%%%%%%%%%%%%%%%%%%%%%%%%%%%%%%%%%
%% § Documento
%%%%%%%%%%%%%%%%%%%%%%%%%%%%%%%%%%%%%%%%%%%%%%%%%%%%%%%%%%%%%%%%%%%%%%%%

% \includeonly{recusion_sobre_numeros_naturales}

% \includexmp{licencia}

\begin{document}
% \dominitoc

\maketitle
\newpage

\noindent
Esta obra está bajo una licencia Reconocimiento--NoComercial--CompartirIgual
2.5 Spain de Creative Commons. 

\vspace*{4ex}

\begin{center}
\fbox{
\begin{minipage}{\textwidth}
\vspace*{1ex}
\noindent                     
{\bf Se permite:}
\begin{itemize}                         
\item copiar, distribuir y comunicar públicamente la obra
\item hacer obras derivadas
\end{itemize}
         
\noindent       
{\bf Bajo las condiciones siguientes:}

\begin{minipage}[t]{20mm}
  \vspace*{0mm}\hspace*{1em}
  \includegraphics[scale=0.35]{deed.png}
\end{minipage}%
\begin{minipage}[t]{14cm}
  \vspace*{0mm}
  {\bf Reconocimiento}. Debe reconocer los créditos de la  obra de la manera
  especificada por el autor.  
\end{minipage}
\\[1ex]

\begin{minipage}[t]{20mm}
  \vspace*{0mm}\hspace*{1em}
  \includegraphics[scale=0.35]{deed-eu.png}
\end{minipage}%
\begin{minipage}[t]{14cm}
  \vspace*{2mm}
  {\bf No comercial}. No puede utilizar esta obra para fines comerciales.
\end{minipage}
\\[1ex]

\begin{minipage}[t]{20mm}
  \vspace*{0mm}\hspace*{1em}
  \includegraphics[scale=0.35]{deed_002.png}
\end{minipage}%
\begin{minipage}[t]{14cm}
  \vspace*{-2mm}
  {\bf Compartir bajo la misma licencia}. Si altera o transforma esta obra, o
  genera una obra derivada, sólo puede distribuir la obra generada bajo una
  licencia idéntica a ésta. 
\end{minipage}

\begin{itemize}    
\item Al reutilizar o distribuir la obra, tiene que dejar bien claro los
  términos de la licencia de esta obra. 
\item Alguna de estas condiciones puede no aplicarse si se obtiene el permiso
 del titular de los derechos de autor.
\end{itemize}
\vspace*{1ex}
\end{minipage}}
\end{center}

\vspace*{4ex}

Esto es un resumen del texto legal (la licencia completa). Para ver una copia
de esta licencia, visite 
\href{http://creativecommons.org/licenses/by-nc-sa/2.5/es/}
     {\url{http://creativecommons.org/licenses/by-nc-sa/2.5/es/}}
o envie una carta a Creative Commons, 559 Nathan Abbott Way, Stanford,
California 94305, USA.

%%% Local Variables: 
%%% mode: latex
%%% TeX-master: t
%%% End: 

\newpage

\newpage

\mbox{} \vspace*{2cm}
\begin{flushright}
\textit{Para Guiomar}
\end{flushright}

\newpage

\tableofcontents
\clearpage

\renewcommand{\chaptername}{Ejercicio}

\chapter*{Introducción}

% \mbox{} \hspace*{1cm} 

\begin{quote}
  ``\textit{The chief goal of my work as an educator and author is to
  help people learn to write beautiful programs.}''

  (Donald Knuth en
  \href{http://www.paulgraham.com/knuth.html}{Computer programming as an art})
\end{quote}

\vspace* {1cm}

Este libro es una recopilación de las soluciones de los ejercicios
propuestos en el blog
\href{https://www.glc.us.es/~jalonso/exercitium}
     {Exercitium}\
     \footnote{\url{https://www.glc.us.es/~jalonso/exercitium}}
durante el curso 2013--14.

El principal objetivo de Exercitium es servir de complemento a la
asignatura de
\href{https://www.cs.us.es/~jalonso/cursos/i1m-13}
     {Informática}\
     \footnote{\url{https://www.cs.us.es/~jalonso/cursos/i1m-13}}
de 1º del Grado en Matemáticas de la Universidad de Sevilla.

Con los problemas de Exercitium, a diferencias de los de las
\href{https://www.cs.us.es/~jalonso/cursos/i1m-13/ejercicios/ejercicios-I1M-2013.pdf}
     {relaciones}\
     \footnote{\url{https://www.cs.us.es/~jalonso/cursos/i1m-13/ejercicios/ejercicios-I1M-2013.pdf}},
se pretende practicar con los conocimientos adquiridos durante todo el
curso, mientras que con las relaciones están orientadas a los nuevos
conocimientos.

Habitualmente de cada ejercicio se muestra distintas soluciones y se
compara sus eficiencias.

La dinámica del blog es la siguiente: cada día, de lunes a viernes, se
propone un ejercicio para que los alumnos escriban distintas soluciones
en los comentarios. Pasado 7 días de la propuesta de cada ejercicio, se
cierra los comentarios y se publica una selección de sus soluciones.

Para conocer la cronología de los temas explicados se puede consultar el
\href{https://www.glc.us.es/~jalonso/vestigium/category/curso/i1m/i1m2013}
     {diario de clase}\
     \footnote{\url{https://www.glc.us.es/~jalonso/vestigium/category/curso/i1m/i1m2013}}.

El código del libro se encuentra en
\href{https://github.com/jaalonso/Exercitium2018}
     {GitHub}\
     \footnote{\url{https://github.com/jaalonso/Exercitium2013}}

\chapter{Iguales al siguiente}
\entrada{Iguales_al_siguiente}

\chapter{Ordenación por el máximo}
\entrada{Ordenados_por_maximo}

\chapter{La bandera tricolor}
\entrada{Bandera_tricolor}

\chapter{Elementos minimales}
\entrada{ElementosMinimales}

% \chapter{Mastermind}
% \entrada{Mastermind}
% 
% \chapter{Primos consecutivos con media capicúa}
% \entrada{Primos_consecutivos_con_media_capicua}
% 
% \chapter{Anagramas}
% \entrada{Anagramas}
% 
% \chapter{Primos equidistantes}
% \entrada{Primos_equidistantes}
% 
% \chapter{Suma si todos los valores son justos}
% \entrada{Suma_si_todos_justos}
% 
% \chapter{Matrices de Toeplitz}
% \entrada{Matriz_Toeplitz}
% 
% \chapter{Máximos locales}
% \entrada{MaximosLocales}
% 
% \chapter{Lista cuadrada}
% \entrada{Lista_cuadrada}
% 
% \chapter{Segmentos de elementos consecutivos}
% \entrada{Segmentos_consecutivos}
% 
% \chapter{Valor de un polinomio mediante vectores}
% \entrada{Valor_de_un_polinomio}
% 
% \chapter{Ramas de un árbol}
% \entrada{Ramas_de_un_arbol}
% 
% \chapter{Alfabeto comenzado en un carácter}
% \entrada{Alfabeto_desde}
% 
% \chapter{Numeración de ternas de naturales}
% \entrada{Numeracion_de_ternas}
% 
% \chapter{Ordenación de estructuras}
% \entrada{Ordenacion_de_estructuras}
% 
% \chapter{Emparejamiento binario}
% \entrada{Emparejamiento_binario}
% 
% \chapter{Ampliación de columnas de una matriz}
% \entrada{Amplia_columnas}
% 
% \chapter{Regiones en el plano}
% \entrada{Regiones}
% 
% \chapter{Elemento más repetido}
% \entrada{Mas_repetido}
% 
% \chapter{Número de pares de elementos adyacentes iguales}
% \entrada{Pares_adyacentes_iguales}
% 
% \chapter{Mayor producto de las ramas de un árbol}
% \entrada{Mayor_producto_de_las_ramas_de_un_arbol}
% 
% \chapter{Biparticiones de una lista}
% \entrada{Biparticiones_de_una_lista}
% 
% \chapter{Trenzado de listas}
% \entrada{Trenza}
% 
% \chapter{Números triangulares con n cifras distintas}
% \entrada{Triangulares_con_cifras}
% 
% \chapter{Enumeración de árboles binarios}
% \entrada{Enumera_arbol}
% 
% \chapter{Elementos con algún vecino menor}
% \entrada{Algun_vecino_menor}
% 
% \chapter{Reiteración de una función}
% \entrada{Reiteracion_de_funciones}
%     
% \chapter{Pim, Pam, Pum y divisibilidad}
% \entrada{PimPamPum}
% 
% \chapter{Código de las alergias}
% \entrada{Alergias}
% 
% \chapter{Índices de valores verdaderos}
% \entrada{Indices_verdaderos}
% 
% \chapter{Descomposiciones triangulares}
% \entrada{Descomposiciones_triangulares}
% 
% \chapter{Número de inversiones}
% \entrada{Numero_de_inversiones}
% 
% \chapter{Sepación por posición}
% \entrada{Separacion_por_posicion}
% 
% \chapter{Emparejamiento de árboles}
% \entrada{EmparejamientoDeArboles}
% 
% \chapter{Eliminación de las ocurrencias unitarias}
% \entrada{Elimina_unitarias}
% 
% \chapter{Ordenada cíclicamente}
% \entrada{Ordenada_ciclicamente}
% 
% \chapter{Órbita prima}
% \entrada{Orbita_prima}
% 
% \chapter{Divisores de un número con final dado}
% \entrada{Divisores_con_final}
% 
% \chapter{Descomposiciones como sumas de n sumandos}
% \entrada{Descomposiciones_con_n_sumandos}
% 
% \chapter{Selección hasta el primero que falla inclusive}
% \entrada{Seleccion_con_fallo}
% 
% \chapter{Buscaminas}
% \entrada{Buscaminas}
% 
% \chapter{Mayor sucesión del problema 3n+1}
% \entrada{Mayor_sucesion_3n%2B1}
% 
% \chapter{Filtro booleano}
% \entrada{Filtro_booleano}
% 
% \chapter{Entero positivo de la cadena}
% \entrada{Entero_positivo_de_la_cadena}
% 
% \chapter{N gramas}
% \entrada{n_gramas}
% 
% \chapter{Sopa de letras}
% \entrada{Sopa_de_letras}
% 
% \chapter{Intercalación de n copias}
% \entrada{Intercala_n_copias}
% 
% \chapter{Eliminación de n elementos}
% \entrada{Elimina_n_elementos}
% 
% \chapter{Límite}
% \entrada{Limite}
% 
% \chapter{Número de palabas que empiezan por mayúscula}
% \entrada{Empiezan_con_mayuscula}
% 
% \chapter{Renombramiento de un árbol}
% \entrada{Renombra_arbol}
% 
% \chapter{Divide si todos son múltiplos}
% \entrada{Divide_si_todos_multiplos}
% 
% \chapter{Ventana deslizante}
% \entrada{Ventana_deslizante}
% 
% \chapter{Representación de Zeckendorf}
% \entrada{Representacion_de_Zeckendorf}
% 
% \chapter{Más cercano cumpliendo una propiedad}
% \entrada{Mas_cercano_cumpliendo_la_propiedad}
% 
% \chapter{Código Morse}
% \entrada{Codigo_Morse}
% 
% \chapter{Producto cartesiano de una familia de conjuntos}
% \entrada{Producto_cartesiano}
% 
% \chapter{Todas tienen par}
% \entrada{Todas_tienen_par}
% 
% \chapter{Sucesiones pucelanas}
% \entrada{Sucesiones_pucelanas}
% 
% \chapter{Producto de matrices como listas de listas}
% \entrada{Producto_de_matrices_como_listas_de_listas}
% 
% \chapter{Inserción en árboles binarios de búsqueda}
% \entrada{Insercion_en_arboles_binarios_de_busqueda}
% 
% \chapter{Matriz permutación}
% \entrada{Matriz_permutacion}
% 
% \chapter{Números con todos sus dígitos primos}
% \entrada{Numeros_con_digitos_primos}
% 
% \chapter{Cadenas de ceros y unos}
% \entrada{Cadenas0y1}
% 
% \chapter{Clausura de un conjunto respecto de una función}
% \entrada{Clausura}
% 
% \chapter{Sustitución en una expresión}
% \entrada{Sustitucion_en_una_expresion}
% 
% \chapter{Laberinto numérico}
% \entrada{Laberinto_numerico}

\end{document}

%%% Local Variables: 
%%% mode: latex
%%% TeX-master: t
%%% End: 

